\chapter{Commentary on \PYS\ 1.24}
\renewcommand{\rightmark}{\MakeUppercase{1.24]}}
\thispagestyle{plain}
\label{trl:on1.24}


``\bht{But\pratikap{p24_1} %\pratikapx %\pratika{atha pradhāna°} 
who is this \Isvara\ who is neither \purusa\ nor \pradhana?}''  Since [\Isvara] is unknown in the discipline of \Samkhya,\index{samkhyasastra@\emph{sāṃkhyaśāstra}} and since proof for the existence of \Isvara\ is necessary, someone who does not understand the meaning of the word ``\Isvara'' asks [the question].%
  \footnote{The reading \emph{sāṃkhyaśāstre prasiddhyabhāvāt, īśvarasadbhāve copapa\emph{tter avaśyatvāt, a}navagateśvara\emph{śabdārthaḥ} pṛcchati} is a result of \inidx{emendation}s.  The reading preserved in L (and its descendants, including the 1952 edition) \emph{sāṃkhyaśāstre prasiddhyabhāvād īśvarasadbhāve copapattim avaśyam anavagateśvaraviśeṣo vā pṛcchati} is hard to construe.  Especially problematic are the two particles \emph{ca} in \emph{copapatti\/°} and \emph{vā\/} in \emph{vā pṛcchati}.  They serve no purpose in L's reading.  In addition, \emph{īśvaraviśeṣa} appears conspicuously close to \purusavisesa\ in the \sutra\ being introduced, but rather unintelligible. Given that the subsequent \sutra s are, according to the author of the YVi, a thesis (\rnYS{1.24}{1.24}) and its proof (\rnYS{1.25}{1.25}), the words \emph{upapatti} and \emph{avaśya} are reasonable to be in the sentence.  Also, the presence of \emph{ca} is hard to explain as a corruption.  Thus I have emended \emph{°papattim a°} to \emph{°papatter a°} and \emph{°vaśyam a°} to \emph{°vaśyatvād a°} (note that in the constituted text I have dissolved sandhi because I insert a comma) to form the second phrase expressing a reason.  For the remaining problem with regard to problematic \emph{°īśvarviśeṣo vā}, I propose to emend it to \emph{anavagateśvaraśabdārthaḥ}.  One consideration here is that the \sutra\ being introduced (1.24) is expected to explain the meaning of the word \Isvara\ (see the commentary on the previous \sutra).  This explanation of the question is prompted largely by the presence of the word \emph{iti} in the \bhasya; it signals that the question comes from a \emph{pūrvapakṣin}. The YVi explains who is asking the question on what ground. A possible explanation of the reading found in the L manuscript is a failed attempt to restore partially illegible text. However, there may be other possibilities.}  %
With regard to this [question] first the [following] \inidx{thesis} (\pratijna) is presented:%
  \footnote{\label{NONTH}The author is consistent in considering that \sutra\ \rnYS{1.24}{1.24} is the thesis and \rnYS{1.25}{1.25} is its proof.  See passages \reftoeditionpar{p24_6}  (translated on p.~\pageref{TSMNI}), %
  \reftoeditionpar{p24_11} (p.~\pageref{TTPR}), %
  and especially \reftoeditionpar{p25_1} (p.~\pageref{YTHPRTJ}).}%

\begin{quote}
\pratikap{p24_2}\st{\Isvara\ is a special \purusa\ who is not tainted by deposits of impurities, deeds (\emph{karman\/}s), and ripening. (\rYS{1.24}{1.24})}
\end{quote}

\bht{Ignorance\pratikap{p24_3} %\pratika{avidyādayaḥ }
(\emph{avidyā}) and so on are}, since they make impure (\emph{kle\-śa\-ya\-nti}), \bht{the impurities (\emph{kleśa\/}s}).  \bht{Deeds (\emph{karman\/}s)} that are performed dependent on them (the impurities) \bht{are good and bad}.  [The word \emph{kuśalākuśalāni} is] an \emph{ekaśeṣa\/} [in which only the last word of those that have the same case ending remained], and thus it represents [all of] ``good [deeds] (\emph{kuśalāni}),'' ``bad [deeds] (\emph{akuśalāni}),'' and ``good and bad [deeds] (\emph{kuśalākuśalāni}).''%
  \footnote{For \emph{ekaśeṣa}, see \rAAP{1.2.64}{1.2.64}: \emph{sarūpāṇām ekaśeṣa ekavibhaktau}\,ff.  The \VMBh\ has a long discussion on that \sutra.}  %
[Of these three groups, the last member,] ``good and bad'' deeds more or less mean mixtures of good and bad.%
  \footnote{This interpretation alludes to \rYS{4.07}{4.7} where yogins' actions are said to be neither white nor black and for others there are three kinds (white, black, and white and black).}  %
\bht{Their result is called ripening}---characterized by the \inidx{birth}, \inidx{length of life}, and \inidx{experience}s.%
  \footnote{This is an allusion to \rYS{2.13}{2.13}: ``When the root exists, there is ripening, [and the ripening is one or a combination of] the birth, length of life, or experiences.''  See p.~\pageref{STR213} for the Sanskrit text.}  %
[The compound] \emph{kleśa\-karma\-vipākā\-śayāḥ\/} is [analyzed as to mean that] they, \inidx{impurities}, and so on, are those that lie around (\emph{āśerata $<$ āśī-}) till \nirvana\labelh{nirvanaone} is reached.  Or, [the compound] \emph{kleśa\-karma\-vipākā\-śayāḥ\/} means heaps of impurities, deeds, and their ripening.\footnote{\label{ASAYASxvt}The second interpretation is based on \Unadisutra\ 4.132: \emph{aśipaṇāyyor uḍāyulukau ca} in which the word \emph{rāśi}\index{rasi@\emph{rāśi}} is said to have derived from the root \emph{aś-}\index{as-@\emph{aś-}} (5th class, ``to pervade'').  By this, the author of the YVi suggests that the word \emph{āśaya}, too, is derived from \emph{aś-} (1st class, ``to pervade'').  This is an alternative to the first interpretation where the root \emph{ā-śī-} was mentioned.  The purpose of presenting two different interpretations of the compound appears to be to de-emphasize the meaning of the word \asaya.  In either interpretation \asaya\ is not referred to as an \inidx{entity}.  Especially in the second interpretation \asaya\ is replaced with \emph{rāśi}, one of many words that signify multitude often without much meaning of their own (such as \emph{jāta, gaṇa}).  Either way, the author understands that \klesa\/s, \karman\/s, and \vipaka\/s stay around and taint \purusa, but there is no residue they leave or a separate storage for them. I did not mark either of the interpretations of the word \emph{āśaya} in the \sutra\ as part of the \bhasya.  All the versions of the \YBh\ \citet{maas:2006} consulted have \emph{tadanuguṇā vāsanā āśayāḥ}, but this version of the \bhasya\ does not reveal how all the words in the compound are related.  I cannot determine whether the author of the YVi intentionally replaced the \bhasya's interpretation with his own or the version of the \bhasya\ he had lacked that sentence.}

\bht{They,\pratikap{p24_4} %\pratika{te manasi }
 although present in the mind,}---since they originate from the activities of the mind---\bht{are attributed to \purusa}.  Why?  \bht{Because [people say] ``He (\purusa) is the one who experiences their result.''}  \bht{In the same way as victory and defeat, although present among soldiers, are attributed to a king}.  The consequence of them (victory and defeat) is bound to the king.

[The\pratikap{p24_5} %\pratika{yo hy anena } 
\Bhasyakara\ continues---]\bht{The one who is not tainted (or affected) by this experience (\emph{bhoga})}\ldots.  [The reference to past] time is not meant [even though a past participle (\emph{aparāmṛṣṭa}) is used]; because [the past participle] is used in a synecdochic sense, i.e., [the past participle \emph{aparāmṛṣṭa} encompasses all the following:] ``[\Isvara] is not tainted, nor will he be tainted, nor has he been tainted.''  ``Tainted (\emph{parāmṛṣṭa})'' means being fit to have a relationship as the locus of the deposits of impurities, deeds, and their ripening.\footnote{I have adopted the reading in L (\emph{kleśa\-karma\-vipākā\-śayā\-nā\-śraya\-sambandha\-yogī}) with as little change as possible: \emph{\emph{kleśa\-karma\-vipākā\-śayā\-nām ā\-śraya\-sambandha\-yogyaḥ}}. The first emendation is to insert another syllable \emph{mā} between \emph{°āsayānā°} and \emph{°śraya°} (hence from \emph{°śayānāśraya°} to \emph{°śayānām āśraya°}). However, the symbol I read as \emph{na} is all but certain. It is dissimilar to ordinary \emph{na} and there appears to be an attempt to make a correction. Another change was from \emph{°yogī} to \emph{°yogyaḥ}. This part, if we look at the transmission, went through changes every time it was copied (from L to M to the edition). It is somewhat understandable since the original reading in L, if left unchanged, contains two sets of almost synonyms (\emph{āśaya} and \emph{āśraya}; \sambandha\ and \yogin. My understanding is that normal \purusa s have the potential of being attached to \emph{kleśa}, etc., as where they take place.}  For, otherwise, the negation (`untainted') would have no effect, since [being tainted by the impurities, etc.,] would not apply [even] to the standard \purusa\ [of \Samkhya], who is free from changes, and has no action.\footnote{M and its descendants, A and \Me, lack this sentence and the subsequent text up to the end of the next paragraph in this translation. This is one complete line (line 3 of folio 21 verso) plus three more symbols at the beginning of the next in L\@.}

Therefore,\pratikap{p24_6} %\pratika{tasmān nitya°}
\labelh{TSMNI} the one who is the mass of \sattva,\footnote{\label{STVNT24}Our author considers the \sattva\ to be the material cause of \Isvara.  This is based on the phrase in the \bhasya\ \emph{prakṛṣṭasattvopādānāt}.  The author understands the meaning of the word \emph{upādāna} as the material cause.  This understanding is made clear in the commentary on the portion of the \bhasya\ that includes the phrase \emph{prakṛṣṭasattvopādānāt}.  See the paragraph \reftoeditionpar{p24_10} that starts with \emph{tatra puruṣa°} translated on p.~\pageref{UPDNMTL25}.  That \Isvara\ has \sattva\ as his material cause and that it is his mind figure prominent point in much of the following discussions on \Isvara's omniscience.} characterised by eternally faultless knowledge and [eternally] unsurpassed sovereignty, \bht{is he, a special \purusa, \Isvara}---this is the thesis [to be proven].

Clarifying\pratikap{p24_7} %\pratika{sūtrārtha°} 
the meaning of the \sutra, [the \Bhasyakara] states: \bht{But then, [emancipated people (\emph{kevalin\/}s),] those who have achieved emancipation (\kaivalya) [are many, as opposed to the one special \purusa, \Isvara]}.  Since it is possible to be untainted by impurities etc., even though it may happen only occasionally, \bht{emancipated people are many [as opposed to the one special \purusa, \Isvara]}.  Are they \Isvara s?   No.  Because there is also the possibility of [their] being tainted.  They are not completely untainted by impurities, etc.  They have both [the possibility of] being tainted (\emph{parāmarśa}) and not being tainted (\emph{aparāmarśa}).  On the other hand, \Isvara\ is only untainted.  For, [when the \sutra\ says] ``untainted'' the word meaning of the two things[, viz., \emph{parāmarśa} and \emph{aparāmarśa},] that should be taken as the characterization [of \Isvara\ to allow the emancipated to be \Isvara s], is not revealed.\footnote{There is a difficulty determining how L reads. The symbol read as a cancelled \emph{t} has somewhat different shape from usual one, and it appears that someone corrected it to \emph{va} or \emph{pa}, but the correction is not inked. This could well be the scribe of M\@. Since \emph{pa} and \emph{va} are often indistinguishable in L, there is a possibility that the symbol read as \emph{vṛ} might have been meant to be a \emph{pṛ}. My conjecture is based on the fact that often \emph{na} and \emph{sa} look similar, and that the symbol that appears to be a cancelled \emph{t} might have been originally meant to be a \emph{m}, and hence \emph{asaṃvṛtaḥ}. The scribe of M intentionally or by mistake read \emph{°pādeyayoḥ} as \emph{°pādeyaḥ}. The dual is preferable (\emph{lectio difficilior}).}  Accordingly,  someone who is completely untainted is expressed by [the word] ``untainted (\emph{aparāmṛṣṭa}).''

\bht{[The emancipated people are many] because they,\pratikap{p24_8} %\pratika{te hi trīṇi } 
having destroyed}, by means of the right view (\emph{samyagdarśana}), \bht{the three [kinds of] bondage}, [viz., the bondage] with regard to \emph{prakṛti}, [the bondage with regard to] \emph{vikṛti}, and [the bondage with regard to] \emph{dakṣiṇā},\footnote{All the commentators (the author of the Yuktidīpikā, the one of the Māṭharavṛtti, the one of the Jayamaṅgalā, and Vācaspati Miśra) on \rSK{44d}{44d} (\emph{viparyayād iṣyate bandhaḥ}) speak of the same three kinds of \emph{bandha} as the author of the YVi.  However, they do not agree on what those \emph{bandha\/}s mean.  See the testimonia on this part of the text.} \bht{have reached emancipation (\kaivalya)}.  \bht{But for \Isvara, the connection with them}, [viz.,] the connection with the impurities and so on, \bht{did not exist, as [it once did] for a liberated person}.  \bht{Because for him (the liberated person), the period of bondage prior [to becoming emancipated] is known [to have existed]} by the mere fact of a state of liberation coming into being.  Because release is [necessarily] preceded by a [state of] bondage.  \bht{Also, [\Isvara] will not have [the connection with them], as someone [still] dissolved in \prakrti\ will.  Because for him}---one who has not yet started the cycle of reincarnation---\bht{it is possible that the future period of bondage will come into existence}.  And it is possible for someone who has started the cycle of reincarnation [but] whose mind is [temporarily] dissolved in \prakrti\footnote{I have emended \emph{cāprakṛtilīnacittasya} to \emph{ca prakṛtilīnacittasya}. As far as the context is concerned, our author should be discussing someone who is temporarily free from the bondages. In that sense, someone whose mind is not dissolved in \prakrti\ is out of place. Considering that \rYBh{1.19}{1.19} (and \rnYBh{1.51}{1.51} to a lesser degree) that our author alludes to by the word \emph{sādhikāra} is about someone who is dissolved in \prakrti\ and experiences as though he is emancipated, this sentence should be about the same kind of people. The idea in \rYBh{1.19}{1.19} is that \emph{citta} completes its responsibility (\adhikara) when one reaches liberation (\moksa/\kaivalya).}---since he is still entitled/obliged [to consume the effects of previous actions]---to have both past and future periods of bondage.

The\pratikap{p24_9} %\pratika{tad idaṃ } 
intention of this whole [\bhasya\ passage] here [from ``But then, those who have achieved emancipation (\emph{kaivalyaprāptās tarhi})'' to ``For this reason, the following was stated: `But He is always liberated [and] always \Isvara' (\emph{etasmād etad bhavati---sa tu sadaiva muktaḥ sadaiveśvara iti})''] is to show that [the past participle] ``untainted (\emph{aparāmṛṣṭa})'' [in the \sutra] is not meant to refer to [the past] time[, as evidenced by the following statement in the \bhasya]---\bht{But\labelh{SATUSD} He is always liberated [and] always \Isvara}.

[The \Bhasyakara\ now]\pratikap{p24_10}%\pratika{tatra puruṣa°}
\labelh{UPDNMTL25} paraphrases the meaning of the next [\sutra], which is already established on the following grounds: \Isvara\ as \purusa\ does not have the quality of sovereignty, because it has been stated in that [preceding \sutra] that \Isvara\ is a special \purusa; furthermore, because the mind \emph{is} endowed with the quality of sovereignty; and finally because an unsurpassed sovereignty implies a connection with an excellent mind. [The paraphrase is the following:] \bht{That eternal excellence of \Isvara\ is the result of having preeminent \sattva\ as the material cause}\labelh{SSVTK24}.  [The word] ``excellence (\utkarsa)'' is [to be understood as] eternal increase of unsurpassed omniscience, sovereignty, and power.

Then\pratikap{p24_11} %\pratika{tatra pratijñā°}
\labelh{TTPR} [the \Bhasyakara] starts [his discussion] preceded by a question, in order to clarify the thesis [that has just been presented, viz., ``That eternal excellence,'' etc.]: \bht{Is there a reason for that [excellence], or is there not?}  The word ``reason (\nimitta)'' is a synonym of ``cause (\kaarana).''  [Thus the question may be paraphrased as] ``Does it have a cause or does it not have a cause?''
% read with arlo by the end of march 2006.

Furthermore,\pratikap{p24_12} %\pratika{kiñ cātaḥ } 
if, first, [the excellence] has a cause, then saying ``always \Isvara'' is inappropriate because [his excellence] will not be eternal.  If, then, [the excellence] does not have a cause, since  excellence is [always] an effect,\footnote{This question is based on the premise that excellence (\utkarsa) is an effect.  However, I am unable to find a relevant passage anywhere that supports the assumption.} there would be the undesirable consequence that it does not exist.  Because people do not know of an effect without a cause.

To\pratikap{p24_13} %\pratika{atrocyate }
\labelh{SSTRISVR24} this [the \Bhasyakara] answers---\bht{It is not that [the excellence] does not have a cause, since its cause is \sastra}---\sastra\ is [\Isvara's] knowledge.\footnote{The discussion here is very brief.  Combined with the textual difficulty (see the next note), this part requires some interpretations.  Especially peculiar is the equation of \sastra\ with knowledge.  With this regard, I find \Sankara's first interpretation of \rBS{1.1.03}{1.1.3} (\emph{śāstrayonitvāt}) in the BSBh shed lights on this part of the YVi.  Note that the \sutra\ is similar to the phrase in the \YBh\ \emph{śāstranimittatvāt} now being explained.  Note also that \Sankara\ gives two alternative interpretations of the \sutra, like the author of the YVi here.  In \Sankara's interpretation of \rBS{1.1.03}{1.1.3} \sastra\ is the \Rgveda, etc.\ (\emph{ṛgvedādi}).  Our author probably understands \sastra\ here in the \YBh\ the same way, viz., all kinds of authoritative teachings.  If he understands the \sastra\ here in a restricted sense, such as the discipline of yoga or any specific discipline, the argument will not carry much force.  Though problematic (see the next note), the compound \emph{sādhyasādhanavyavahitaviprakṛṣṭa°} obviously refers to all kinds of \sastra.

\Sankara\ understands \rBS{1.1.03}{1.1.3} \emph{śāstrayonitvāt} as ``Because [\Brahman{}] is the source of \sastra.''  This is different from the interpretation here.  The author of the YVi reads \sastra\ as the cause.  However, \Sankara's first interpretation makes it so that the existence of \sastra\ is the evidence of the existence of an all-knowing personality (\sarvajna, \purusavisesa).  He reasons that any \sastra\ presupposes an intelligence greater than the information given in the \sastra\ (see note \ref{BSBH113} on p.~\pageref{BSBH113}).  Thus he sees \sastra\ (the Ṛgveda, etc.)\ as the manifestation of the knowledge \Brahman{} has.  Again, assuming a similar idea behind the equation of \sastra\ and knowledge helps us comprehend our text.  In the equation knowledge (\jnana) probably means the knowledge \Isvara\ has.}  [To paraphrase: his excellence is not without a cause] because \sastra\ is the cause.  For, [his] knowledge, rooted in [the whole] of goals and their means, including things that are concealed or remote [as evidenced in the Vedas and other \sastra\/s], has everything as its object, and is, like the nature of material, eternal.\footnote{I have emended \emph{dravyasvabhāvasādhyasādhana°} in the manuscripts to \emph{dravya\-sva\-bhāva\-vat sādhyasādhana°}.  Main reason is that the author of the YVi never uses the word (\emph{dravya})\emph{svabhāva}/\emph{svarūpa} to refer to the object of knowledge whereas the rest of the compound is apparently about the object of knowledge.  Rather, when the author discusses the knowledge \Isvara\ has (\reftoeditionpar{p25_33} \emph{etena\ldots}, trl.\ p.~\pageref{ETEN}\,ff.; \reftoeditionpar{p25_141} \emph{atha cakṣurādi°\ldots}, trl.\ on p.~\pageref{ATHACKS}\,ff., etc.), he emphasizes that it is the nature of (\citta)\sattva\ that it is connected to everything and takes the shape of the object.  He regularly uses the similes of light (\prakasa) or the space (\akasa) in those discussions.  (A similar idea is discussed in \rBSBh{1.1.05}{1.1.5}.)  He also defends that knowledge of \Isvara\ may be natural (in the paragraph \emph{kiñ ca\ldots}\ below), in addition to following the \bhasya\ in stating that the \sastra/\jnana\ has a cause.  It seems natural that he would have introduced the idea of innate knowledge already here.  %[commented out on 20100421] It is notable that the compound \emph{dravyasvabhāva} does not occur frequently.
  The expression \dravyasvabhava\ does not appear again in the YVi.  \Sankara\ uses it once in \rBSBh{2.1.24}{2.1.24} to refer to the natural process of milk turning into yoghurt or water turning into snow or hail (\emph{yataḥ kṣīravad dravyasvabhāvaviśeṣād upa\-padya\-te\dd\ yathā hi loke kṣīraṃ jalaṃ vā svayam eva dadhihimakarakādibhāvena pariṇamate ’nāpekṣya bāhyaṃ sādhanam, tathehāpi bhaviṣyati\dd}).  This is part of the answer to the question whether \Brahman{} needs external means to be the creator.  The argument (``because it so happens'') is similar to the one in the YVi where the author argues that \Isvara\ does not need a means to know everything.  % [commented out on 20100421] Other usages I am aware of are once in the \VP\ (Sādhanasamuddeśa 138) and once in \NBh\ 3.1.49.% [commented out on an unknown date] A similar notion that the nature of something is eternal is expressed in the YVi in the discourse against the view that the origin of the universe is a natural process (see pp.~\pageref{NTRLVD}ff.).  Such being the case, I consider that the word \emph{dravyasvabhāva} should be closely associated with the eternal knowledge whose object is everything (\emph{sarvaviṣayaṃ jñānan nityam}) in the sentence.  Comparing \Isvara's knowledge to the nature of a substance as something innate to him appears sensible.  

I understand the rest, \emph{sādhyasādhanavyavahitaviprakṛṣṭāspadam}, to refer to the content of \Isvara's knowledge or \sastra.  All the uses of the compound \emph{sādhyasādhana} in the YVi are related to \sastra\ (see paragraphs \emph{sambandho 'pi\ldots} \reftoeditionpar{p01_9} translated on p.~\pageref{SMBDHPI}, \emph{yogānuśāsanam\ldots} \reftoeditionpar{p01_27} translated on p.~\pageref{YGNSSM}, \emph{agnihotrādi°\ldots} \reftoeditionpar{p25_26} translated on p.~\pageref{AGNIHTR}), and the relationship between the two is one of the matters that one expects to be discussed in a \sastra.  In addition, the author has the notion that everything in the universe is connected in the form of means and goals (and their relationships are taught in \sastra).  See the syllogism \emph{sarvam etat\ldots}\ \reftoeditionpar{p25_23} translated on p.~\pageref{UNIGOAL}.  Expressions based on the same idea using the compound \emph{sādhyasādhana} are found several times in the \BAUBh:  \rnBAUBh{1.4.06}{1.4.6} \emph{sādhyasādhanalakṣaṇasya vyākṛtasya jagataḥ} \citep[66]{bau}, \rnBAUBh{1.4.07}{1.4.7} \emph{idam iti vyākṛtanāmarūpātmakaṃ sādhyasādhanalakṣaṇaṃ yathāvarṇitam abhidhīyate} \citep[67]{bau}, \rnBAUBh{1.5.02}{1.5.2} \emph{sarvo lokaḥ sādhyasādhanalakṣaṇaḥ} \citep[122]{bau}, \rnBAUBh{2.1.01}{2.1.1} \emph{sādhyasādhanalakṣaṇaṃ vyākṛtaṃ jagat} \citep[138]{bau}.  And \emph{vyavahitaviprakṛṣṭa} is more or less a standard expression (usually with \emph{sūkṣma}) to refer to objects that are beyond the reach of sense faculties.

There appear to be several assumptions and steps of logic in this short explanation of the \bhasya.  First, it is assumed that \sastra\ originates from \Isvara, or at least \sastra\ is a manifestation of \Isvara's knowledge.  Otherwise, the author would not paraphrase \sastra\ with \jnana.  (See the previous note.)  And hence \Isvara's knowledge is equal to or greater than \sastra.  Another assumption is that \sastra\ is about everything.   And therefore \Isvara's knowledge covers everything and is eternal.  Accordingly, this knowledge gives eternal advantage or excellence to \Isvara.  Yet another assumption here is that the greater knowledge results in superiority.  It is notable that the author appears to be using the existence of \sastra\ as the evidence of \Isvara's omniscience.  Cf.\ note \ref{BSBH113} on p.~\pageref{BSBH113}.}

If\pratikap{p24_14} %\pratika{yady evaṃ } 
it is so, then, since knowledge is [always] preceded by learning something from somewhere else, there would be no excellence prior [to that learning].  Considering this possibility, [the \Bhasyakara] says \bht{What does \sastra\ then have as its cause?}  If [\Isvara's] knowledge is innate, then [this knowledge] would be the cause of lack of excellence [rather than that of excellence], like that of a drunkard or an insane person [whose knowledge comes from nowhere].  But if, on the other hand, it (\sastra, i.e., knowledge) does have a cause, then it will be thought that there is no excellence [of \Isvara] prior to that cause.

First,\pratikap{p24_15} %\pratika{na tāvan } 
it is not that it (\sastra, i.e., \Isvara's knowledge) does not have a cause.  Because \bht{[it] has [\Isvara's] preeminent \sattva\ as its cause}.  This means that [the \sastra] relies on [his] preeminent \sattva.  For this very reason, there is no fault in knowledge  being innate.

Furthermore,\pratikap{p24_16} %\pratika{kiñ ca } 
even if [his] knowledge is innate, this will not be the cause of inferiority [of \Isvara], like the knowledge of drunkards, etc., [is].  Because it relies on \sattva\ eternally detached from impurities (\emph{kleśa}), etc.  As to this: knowledge, impressions, memory and recollection\footnote{I have emended the reading \emph{jñāna\-saṃskāra\-smṛti\-pra\-ti\-bandhānām} found in the manu\-script L and its descendants (no reading from T is available) to \emph{jñāna\-saṃskāra\-smṛti\-prati\-sandhānānām}, following the usage in YVi 1.4, for example.  The word \emph{pratibandha} in this compound does not fit. The words \emph{smṛti} and \emph{pratisandhāna} form a standard combination in many other texts.} have a seed-and-sprout-like connection, without [clear] starting point, due to each other's being cause and the effect.  In precisely the same manner, \dag there is a beginning-less and endless connection between the \sastra\ and excellence, which are active in the mind of \Isvara, in the form of an eternally active process.\dag\footnote{I have marked this part of the text as unrepairable. This is not so much as the text is incomprehensible, but in relation to the root text being commented. Having much in common with the text of the \bhasya\ [see \citet[37]{maas:2006}], this part may be meant to embed the words of the \bhasya, or it may simply be a close paraphrase. I expect the former since that is our author's style. There are several problems. Maas's edition, based on the manuscript tradition of the \YBh, reads: \emph{etayoḥ śāstrotkarṣayor īśvarasattve vartamānayor anādiḥ sambandhaḥ}. However, our text here does not have \emph{etayoh}; it has \emph{śāstraprakaṣayoḥ} instead of \emph{śāstrotkarṣayoḥ}; it has \emph{īśvaracetasi} instead of \emph{īśvarasattve}; it has \emph{pravartamānayoḥ} instead of \emph{vartamānayoḥ}; and it has \emph{anādyantaḥ} instead of \emph{anādiḥ}. First, we cannot know whether \emph{etayoḥ} was originally in our text or dropped during the transmission. Our author knew the word in the \bhasya\ (see the next paragraph). If this part was meant to be a gloss embedding all the words of the \bhasya, then the lack of \emph{etayoḥ} is an error. If, on the other hand, if this is meant to be a close paraphrase, not necessarily so. As for \emph{śāstraprakarṣayoḥ}, this is clearly an error; it should be \emph{śāstrotkarṣayoḥ} as in Maas's. That in fact is how our author reads, as seen in the next paragraph. As for \emph{īśvaracetasi}, we could theorize this is a case of eye-skip from the original reading \emph{īśvarasattva īśvaracetasi}. Glossing the word \emph{īśvarasattve} with \emph{īśvaracetasi} is reasonable, for the word is ambiguous (what does \emph{sattva} mean here?).  Nonetheless, our author did know the reading \emph{īśvarasattve}, as seen in the opposition two paragraphs below. Similarly, the reading \emph{pravartamānayoḥ} is suspicious of being an error since our author knew the reading \emph{vartamāna°} (again, see two paragraphs below). As for \emph{anādyantaḥ} against \emph{anādiḥ}, we at least know that our author did not read \emph{anādiḥ} as in Maas's reconstruction of the \bhasya. He read \emph{anādyantaḥ} or \emph{nirādyantaḥ} as found in the next paragraph. Alongside the next paragraph where an alternative interpretation is given, we can be more-or-less certain that our author knew the \bhasya\ that read similarly to other manuscript tradition of the \YBh\ with the difference of \emph{anādyantaḥ} or \emph{nirādyantaḥ}. Considering all these, I could propose the reading of this part as \emph{tathai\textbf{vaitayoḥ śāstrotkarṣayor īśvarasattva} īśvaracetasi nityapravṛttaprabandharūpeṇa \textbf{vartamānayor anādyantaḥ sambandhaḥ}}. However, this is only speculative and I keep the text found in our manuscripts in the edition. Note that in the following paragraph, which proposes an alternative interpretation of the same \bhasya\ passage, too, not all the words of the \bhasya\ the author knew were embedded. This may be an indication that our author did not thoroughly follow the principle to embed all the words in the root text in the exact order found there.}  In that [series of causes and effects], it is the excellence which is the effect of knowledge; the knowledge, too, is its effect.

An\pratikap{p24_17} %\pratika{anyeṣāṃ }
\labelh{ANYSMV24} interpretation by others:\footnote{\label{ANYSM24NOTE}Cf.\ this interpretation with the alternative interpretation of \BS\ 1.1.3 that starts with \emph{atha vā}.  There \Sankara\ interprets the \sutra\ as telling that \sastra\ informs of \Brahman, just like the interpretation here.} the word ``reason (\emph{nimitta})'' expresses a means of knowledge (\pramana).  \emph{Śāstra\/} is the reason, i.e., the means of knowledge, of that [excellence of \Isvara].  Because his excellence is proven by it.  What then proves [the reliability] of \sastra?  The spotless \sattva\ of \Isvara\ is the proof.  For the reliability of \sastra\ results from the fact that it has been propounded by someone with pure \sattva, as in cases such as the \Manava [dharmaśāstra].  So the \sruti\ says, ``Anything Manu said was a medicine.''\footnote{Note that the reading of the \Taittiriyasamhita\ quote is not the same as found in its editions. (See p.~\pageref{TS22102BS211} for the reading.) The difference is probably due rather to changes in the transmission of our text (more classical than Vedic) than to a different version of the \Taittiriyasamhita\ known to our author. Yet I did not adjust the reading for the small possibility that our author indeed knew the reading found here.}  Similarly, people say, ``The master has said [such and such and therefore it must be true].''  \bht{There is a beginningless and endless relationship between these two, \sastra\ and excellence,} as the means of knowledge and the object to be known.%\footnote{Note that the author of the YVi continues the commentary based on the interpretation that he has said to belong to others.}

[Objection]\pratikap{p24_18} %\pratika{nanu ca } 
It is not at all possible that his excellence is proven by a \sastra\ that is present in \Isvara's [own] \sattva\ [and not external to it].\footnote{The suggestion by the editors of the 1952 edition to read \emph{vartamāna°} as \emph{avartamāna°} is possible, because the manuscripts seldom record \avagraha.  However, in order for the objection to have any effect, both parties (the opponent and the proponent) should accept that \sastra\ does not exist in \Isvara's \sattva.  But such an assumption would not be accepted by the proponent who has just said that \sastra\ exists in \Isvara's \sattva.   If this is an argument on the authenticity of \sabda/\agama, there is another implicit argument we must assume in reading \emph{avartamāna}: a \sastra\ existing outside of the \sattva\ of the speaker is not authentic.  But such an argument is not found in any text.  Rather, the question from the opponent makes sense if the question is based on the assumption that as long as a \sastra\ resides only in \Isvara's \sattva, it will be invisible and will not serve to inform us of \Isvara's superiority.}

[Answer]\pratikap{p24_19} %\pratika{naiṣa doṣaḥ } 
There is no such fault.  Since [\sastra] originates from that [\sattva\ of \Isvara], it [is said to] exist there. \emph{Śāstra\/} indeed exists there while it also originates there, for the very reason that [\Isvara] is omniscient.\footnote{I have some doubt about the reading \emph{sarvajñatvāt} found in L (and its descendants, M, A, and \Me).  As it stands, presumably the argument is that since \Isvara\ is omniscient, he still knows what flowed out as \sastra, i.e., producing/emitting \sastra\ does not mean the loss of what he knew.  It does not, however, explain why the \sastra\ that originates from \Isvara's \sattva\ can still be in that \sattva.  An emendation to \emph{sarvagatatvāt} might be possible.  The argument makes more sense if it answers why this \sastra\ that originates from \Isvara's \sattva\ can be said to exist there.}  Also in our daily life [we find] that when A originates from B, A exists in B, for example: the cloth and so on exists in the thread and so on.\footnote{Note that this is a typical explanation of \satkaryavada.} That [\sastra] originates there is known from inference (\anumana) and from reliable testimony (\agama).  The excellence [of \Isvara] is proven by \sastra; the reliability of \sastra\/ is due to \Isvara.  Accordingly, since they depend on different factors, one [can]not [object] that the excellence and \sastra\/ depend on each other [circularly].

[Objection]\pratikap{p24_20} %\pratika{nanu ca } 
But if the authority of \Isvara\ is due to \sastra, and if the authority of \sastra\/ is due to his authority, there would be [the fault of circular] mutual dependence.

Answer:\pratikap{p24_21} %\pratika{ucyate } 
since the authority of \Isvara\ is established by inference, there is no such fault.\labelh{AUTHINF}

\bht{For\pratikap{p24_22} %\pratika{etasmād etad } 
this reason, the following has been stated [previously].}\footnote{See page~\pageref{SATUSD}.} For what reason?  Because the unsurpassed knowledge and excellence that rely on \Isvara's preeminent \sattva\/ are eternal, as the process in the form of cause and effect is eternal. % this is some sort of eternal engine idea.
[Therefore the following was stated:] \bht{``But he is always liberated [and] always \Isvara.''}

About this\pratikap{p24_23} %\pratika{atrāha }
\labelh{ATISAYAM24} [the \Bhasyakara] says: \bht{And thus, his sovereignty is without an equal or superior}.  This is the conclusion (\emph{phala}) of the proof (\pramana) about to be stated (viz., inference) or the approximate meaning of this [following] \sutra.

[The\pratikap{p24_24} %\pratika{atiśayavimuktatāṃ } 
\Bhasyakara] elaborates what it is [for his sovereignty] to have no superior: \bht{In the first place, it is not surpassed by another sovereignty.}  Why?  \bht{The superior sovereignty} you think of \bht{is exactly the sovereignty} of his that I have spoken of.  This means that if a sovereignty which surpasses another sovereignty is present anywhere, that exactly is \Isvara.
% up to here is read with Arlo on 4/4/2006.

\bht{Nor\pratikap{p24_25} %\pratika{na ca tat°} 
is there sovereignty that equals his}; because equal eminence of properties is impossible.  For, it is impossible for a kingdom to have two kings, [nor is it possible] for one king to have two kingdoms.  [The \Bhasyakara] substantiates: \bht{[If it were possible that there is a sovereignty that equals his,] then the undesirable consequence that the two [sovereignties] are equal would follow}.  [This is undesirable] because it is impossible for either of the two [sovereignties] to act on the other [sovereignty] without annulling the other's volition[, and this amounts to neither of them being sovereignty].\footnote{These two paragraphs pose difficulties.  Difficulties include the fact that the \bhasya\ our author is commenting appears to be quite different from any other versions of it, as well as the fact that there are signs of corruptions in the text.}

\bht{Nor\pratikap{p24_25a} is there only effect on two objects simultaneously} with one volition.  \bht{Because} when [one] object is desired, \bht{this precludes another object} having the same eminence.  If A has the same eminence as B, then A is annulled by that B\@.  \begin{uncertain}And this [possibility of A having the same eminence] is annulled by that [B]\@.\end{uncertain}%
\footnote{The reading \emph{te ca na vācā bādhanam} is not comprehensible.  The repeated appearance of \emph{c} and \emph{v}, which look similar in the Malayalam script, makes one suspect a corruption.  The conjecture is based on the understanding that annulment should be done on the eminence, not in the sense one object disappears.  Still, it is not very clear how this sentence relates to the next quotation.%  This emendation was suggested by Arlo Griffiths.
}  Elsewhere we will state: ``[In manifestations of three \guna s, one \guna's] dominances in a characteristic, and [one \guna's] dominances in a function, do not coexist [with the others' dominance].  But [in all manifestations] commonalities of [the three \guna s, either dominant or subsidiary] accompany the dominant ones.''%
\footnote{While the argument of the \bhasya\ interpreted here is more or less straightforward, I fail to see the exact purpose of the additional information given in the commentary.  It does not appear to supply much support to the argument on the impossibility of the existence of two \aisvarya s of exactly the same value.  The author rather appears to stress that no two things can be equal.%

	The author of the \YD\ ascribes the fragment quoted in \rYBh{2.15}{2.15} (referred to again in 3.13) to \Varsaganya.  The author of our text ascribes this fragment to a \emph{tantra} \citep[167]{yvi}.  We may naturally assume that this ``Tantra'' is the \Sastitantra, variously ascribed to \Kapila, \Pancasikha, and \Varsaganya\ \citep[127]{eoip:samkhya}.}%

\bht{Therefore,\pratikap{p24_26} %\pratika{tasmād yasya } 
the one to whom sovereignty, that lacks any equal or superior, belongs is \Isvara.}  Thus \Isvara, who is neither \pradhana\/ nor \purusa, a special \purusa, has been proved.

